\chapter{Área Básica de Ingresso e Trajetórias Formativas}
\label{chp:abi}

Este capítulo detalha a proposta de organização do \ppccurso em uma Área Básica de Ingresso (ABI) compartilhada com o curso de Engenharia de Automação, Robótica e Inteligência Artificial. Trata-se de dois cursos de graduação distintos — com PPC, matriz curricular, perfil de egresso, requisitos de integralização, colação de grau e diploma próprios —, cujas matrizes possuem uma etapa formativa inicial comum. A implantação do ingresso unificado, da escolha de trajetória e de qualquer mecanismo de continuidade entre os cursos depende de ato institucional específico e de disposições simétricas nos dois PPCs.

A origem institucional desta ABI é comum aos dois cursos: ambos foram contemplados pelo Programa Universidades Inovadoras e Sustentáveis~\cite{programa_universidades_inovadoras_sustentaveis}, que concedeu à FEELT novas vagas docentes e de técnicos de laboratório — no caso deste curso, sustentando sua criação como curso inteiramente novo; no caso da Engenharia de Automação, Robótica e Inteligência Artificial, sustentando sua reformulação e renomeação a partir do curso já existente de Engenharia de Controle e Automação, com mudança para o turno noturno e ampliação de 15 para 30 vagas ofertadas (\autoref{sec:historico}). Essa origem comum é o que justifica institucionalmente a organização dos dois cursos em uma mesma ABI, e não apenas a afinidade técnica entre as áreas.

O conceito que orienta esta organização curricular pode ser sintetizado da seguinte forma: o \ppccurso constitui uma \textbf{formação tecnológica completa, autônoma e terminal}; o curso de Engenharia de Automação, Robótica e Inteligência Artificial constitui essa mesma \textbf{formação tecnológica comum, ampliada} por fundamentação científica, matemática e metodológica própria da Engenharia, que sustenta competências de concepção, modelagem, análise, projeto, dimensionamento, integração e validação de sistemas de maior abrangência e complexidade (\autoref{sec:distincao_perfis}). A base tecnológica comum não torna os dois cursos equivalentes, nem reduz o \ppccurso a uma etapa incompleta da Engenharia: trata-se de duas formações superiores distintas, cada uma legítima e completa segundo seus próprios objetivos, terminalidades e perfis profissionais de egresso.

\section{Caracterização da Área Básica de Ingresso}

A Área Básica de Ingresso (ABI) é uma forma de organização acadêmica em que o estudante ingressa em uma etapa inicial comum associada a mais de um curso e, posteriormente, segue uma das graduações vinculadas. A ABI não é curso autônomo, não confere titulação própria e não substitui a identidade regulatória dos cursos. As Normas Gerais da Graduação da UFU~\cite{ufu_congrad_46_2022} não regulamentam, por si só, todos os elementos operacionais desta proposta entre dois cursos distintos; por isso, a forma de ingresso, o registro acadêmico e a escolha de trajetória deverão constar de ato institucional próprio.

O compartilhamento do núcleo tecnológico inicial entre os dois cursos decorre da afinidade entre as áreas de Automação Industrial e de Controle e Automação, da racionalidade acadêmica na integração curricular entre cursos correlatos, do aproveitamento de conhecimentos comuns e do propósito institucional de constituir, para o(a) estudante, uma trajetória formativa integrada e progressiva — sem que isso implique fundir os dois cursos em um único perfil de egresso ou tratar um deles como decorrência incompleta do outro.

A escolha da trajetória ao final da etapa comum (\autoref{sec:escolha}) não deve ser interpretada como seleção entre um curso completo e um curso incompleto, tampouco como decisão entre concluir integralmente uma formação ou apenas parte dela: trata-se da escolha entre duas formações superiores diferentes, ambas legítimas, ambas completas segundo seus próprios objetivos formativos, e ambas com terminalidade profissional própria (\autoref{sec:distincao_perfis}).

A ABI não deve ser confundida com:

\begin{itemize}
    \item ciclo básico ou núcleo comum genérico, sem organização formal de ingresso e de escolha entre graduações;
    \item migração curricular — mudança de estudantes vinculados a uma matriz curricular anterior para uma nova matriz do mesmo curso (\autoref{sec:migracao_curricular});
    \item transferência de curso ou novo ingresso por processo seletivo;
    \item mudança automática de curso; e
    \item dupla diplomação automática — a obtenção de cada diploma depende do cumprimento integral dos requisitos da respectiva graduação (\autoref{sec:diplomas}).
\end{itemize}

\textbf{Pendência institucional crítica:} aprovar o ato de criação e a denominação oficial da ABI, com cursos e códigos vinculados, vagas, ingresso, registro acadêmico, procedimento de escolha, autoridade decisória e regras de transição.

\section{Distinção entre as formações e os perfis profissionais}
\label{sec:distincao_perfis}

Embora o \ppccurso e o curso de Engenharia de Automação, Robótica e Inteligência Artificial compartilhem um núcleo tecnológico comum, as duas formações possuem objetivos, competências, terminalidades e perfis profissionais distintos. O \ppccurso constitui uma graduação tecnológica completa e autônoma, orientada à aplicação, à integração, à implementação, à operação, à supervisão, à manutenção, à gestão e ao desenvolvimento de sistemas de automação industrial (\autoref{chp:perfil_egresso}). O curso de Engenharia de Automação, Robótica e Inteligência Artificial incorpora essa base tecnológica e a amplia por meio de formação científica, matemática, metodológica e multidisciplinar própria da Engenharia, desenvolvendo competências para conceber, modelar, analisar, dimensionar, projetar, integrar e validar sistemas e soluções de engenharia de maior abrangência e complexidade.

A ampliação promovida pelo curso de Engenharia de Automação, Robótica e Inteligência Artificial não se resume a um acréscimo de carga horária, de disciplinas ou de conteúdos matemáticos: a fundamentação científica e matemática adicional sustenta o desenvolvimento de capacidades próprias de engenharia — formulação e análise de problemas complexos, modelagem matemática de sistemas físicos e dinâmicos, análise de estabilidade e de desempenho, concepção de arquiteturas de controle e automação, dimensionamento e projeto de sistemas integrados, simulação e validação de soluções, otimização de processos e recursos, análise de riscos e tomada de decisão técnica considerando aspectos econômicos, sociais, ambientais e éticos, entre outras. O(a) engenheiro(a) formado(a) por esse curso domina também os conteúdos tecnológicos desenvolvidos no \ppccurso, mas os aplica em um contexto de maior profundidade científica, capacidade de abstração, integração multidisciplinar e responsabilidade técnica sobre sistemas complexos.

Essa distinção é qualitativa, não apenas quantitativa, e pode ser sintetizada como segue:

\begin{itemize}
    \item \textbf{Perfil predominante do(a) \ppctitulacao:} formação especializada, tecnológica e aplicada, direcionada à integração, implementação, operação, supervisão, manutenção, gestão e desenvolvimento de soluções de automação industrial em contextos delimitados e compatíveis com essa formação.
    \item \textbf{Perfil predominante do(a) engenheiro(a) formado(a) pelo curso de Engenharia de Automação, Robótica e Inteligência Artificial:} formação ampla, científica, matemática, tecnológica e multidisciplinar, direcionada à análise, modelagem, concepção, dimensionamento, projeto, integração, validação e gestão de sistemas e soluções de engenharia de maior abrangência e complexidade.
\end{itemize}

Essa distinção não é absoluta: ambos os perfis podem participar de atividades de projeto, desenvolvimento e integração de sistemas de automação, porém com níveis de profundidade, abrangência, complexidade e responsabilidade técnica coerentes com a respectiva formação — não sendo adequado descrever o(a) \ppctitulacao como profissional que apenas executa ou opera, nem o(a) engenheiro(a) como o(a) único(a) capacitado(a) a projetar.

O quadro a seguir sintetiza essa distinção:

\begin{longtblr}[
    theme = ppc,
    caption = {Distinção entre o \ppccurso e o curso de Engenharia de Automação, Robótica e Inteligência Artificial},
    label = {tab:comparativo_cst_engenharia},
]{
  colspec = {Q[l,m,wd=30mm]X[l,m]X[l,m]},
  rowhead = 1,
  rowsep = 1pt,
  row{odd} = {bg=CinzaClaro},
  row{1} = {bg=AzulEscuro, fg=white},
}
\textbf{Dimensão} & \textbf{\ppccurso} & \textbf{Engenharia de Automação, Robótica e Inteligência Artificial} \\
Natureza & Graduação tecnológica & Bacharelado em Engenharia \\
Terminalidade & Completa e autônoma & Completa e autônoma \\
Ênfase & Aplicação, integração e gestão de tecnologias de automação & Concepção, modelagem, projeto e validação de sistemas de engenharia \\
Formação científica e matemática & Aplicada às necessidades do campo tecnológico & Ampla e aprofundada, para análise e projeto de sistemas complexos \\
Complexidade predominante & Sistemas e processos tecnológicos delimitados & Sistemas multidisciplinares de maior abrangência e complexidade \\
Perfil do egresso & \ppctitulacao & Engenheiro(a) em Automação, Robótica e Inteligência Artificial \\
Diploma & Diploma próprio de graduação tecnológica & Diploma próprio de bacharelado em Engenharia \\
Base comum & Núcleo tecnológico integrado (\autoref{sec:etapa_comum}) & Núcleo tecnológico integrado, acrescido da formação própria de Engenharia \\
\end{longtblr}

\section{Cursos integrantes}

Integram esta Área Básica de Ingresso (ABI):

\begin{enumerate}
    \item o \ppccurso, objeto deste PPC, de grau \ppcgrau, com duração mínima de \ppctempominimo, ofertado pela Faculdade de Engenharia Elétrica (FEELT) da UFU; e
    \item a Engenharia de Automação, Robótica e Inteligência Artificial, curso de graduação de grau Bacharel, com duração mínima de dez períodos (cinco anos), objeto de PPC próprio, elaborado separadamente do presente documento.
\end{enumerate}

Os dois cursos mantêm bases regulatórias distintas — o \ppccurso está estruturado em conformidade com o Catálogo Nacional de Cursos Superiores de Tecnologia (CNCST, 4ª edição)~\cite{mec_cncst_2024} e com as Diretrizes Curriculares Nacionais Gerais para a Educação Profissional e Tecnológica~\cite{cne_cp_1_2021} (\autoref{chp:identificacao}; \autoref{chp:justificativa}), ao passo que a Engenharia de Automação, Robótica e Inteligência Artificial se estrutura conforme as Diretrizes Curriculares Nacionais do curso de graduação em Engenharia. A articulação pela ABI não uniformiza essas bases regulatórias: cada curso preserva seu próprio marco normativo, sua própria matriz curricular e seu próprio perfil de egresso.

\textcolor{red}{[CONFIRMAR A FORMA DE INGRESSO E O CÓDIGO ACADÊMICO INICIAL]}

\textcolor{red}{[CONFIRMAR A EXISTÊNCIA E A DISTRIBUIÇÃO DE VAGAS]}

\section{Etapa formativa comum}
\label{sec:etapa_comum}

Os cinco primeiros períodos da trajetória acadêmica na ABI são comuns entre o \ppccurso e a Engenharia de Automação, Robótica e Inteligência Artificial. É importante destacar que esses cinco períodos \textbf{não são adicionais} à duração de nenhum dos dois cursos:

\begin{itemize}
    \item correspondem aos cinco primeiros períodos da própria matriz curricular do \ppccurso, cuja duração mínima total é de seis períodos (\ppctempominimo; \autoref{sec:regime}); e
    \item correspondem, simetricamente, aos cinco primeiros períodos da matriz curricular da Engenharia de Automação, Robótica e Inteligência Artificial, cuja duração mínima total é de dez períodos.
\end{itemize}

Ou seja: um mesmo bloco de cinco períodos é, ao mesmo tempo, o início da trajetória do \ppccurso e o início da trajetória da Engenharia de Automação, Robótica e Inteligência Artificial — não um acréscimo prévio a nenhuma das duas. A organização geral da ABI é, portanto:

\begin{center}
Ingresso na Área Básica de Ingresso (ABI) \\
$\downarrow$ \\
realização dos cinco primeiros períodos comuns \\
$\downarrow$ \\
escolha da graduação que o estudante pretende concluir inicialmente (\autoref{sec:escolha}) \\
$\downarrow$ \\
continuidade na trajetória específica escolhida (\autoref{sec:trajetoria_tecnologo}; \autoref{sec:trajetoria_engenharia})
\end{center}

Verificação realizada nesta revisão: os componentes curriculares ativos dos cinco primeiros períodos são atualmente idênticos, código a código, entre a matriz curricular deste curso e a da Engenharia de Automação, Robótica e Inteligência Artificial — confirmando que a etapa comum é, de fato, integralmente compartilhada.

\subsection{Módulos ACE comuns e específicos}

Essa identidade curricular da etapa comum se estende às Atividades Curriculares de Extensão (ACE), organizadas em ambos os cursos em módulos curriculares de extensão (\autoref{sec:ace} de cada PPC). Os módulos \textbf{ACE I} (4º período, 90h), \textbf{ACE II} (5º período, 90h) e \textbf{ACE III} (6º período, 60h) são, código a código e hora a hora, o mesmo módulo curricular nos dois cursos, o que permite seu aproveitamento integral entre as duas matrizes em caso de permanência de vínculo (\autoref{sec:permanencia}), sem repetição de atividades extensionistas. O módulo ACE III é comum mesmo ofertado no 6º período — já o primeiro período em que as demais disciplinas das duas matrizes se distinguem (Seção~\ref{sec:etapa_comum}) —, o que faz deste \ppccurso um curso sem nenhum módulo ACE específico próprio: seus três módulos são integralmente compartilhados com a Engenharia. Os módulos \textbf{ACE IV} e \textbf{ACE V} (7º e 8º período, 60h cada), sem correspondência na matriz deste curso, de seis períodos, são específicos da Engenharia de Automação, Robótica e Inteligência Artificial, com eixos temáticos de aprofundamento em robótica, inteligência artificial e pesquisa aplicada (\autoref{sec:distincao_perfis}). O detalhamento de cada módulo — identificação, objetivos, resultados de aprendizagem, formas de integralização, avaliação, registro e governança — consta da Seção~\ref{sec:ace} de cada PPC.

\section{Escolha ao final do quinto período}
\label{sec:escolha}

Ao completar o quinto período — 2 (dois) anos e 6 (seis) meses de curso —, o estudante deverá formalizar a escolha da graduação que pretende concluir inicialmente: o \ppccurso ou a Engenharia de Automação, Robótica e Inteligência Artificial. Essa escolha define a trajetória específica que o estudante seguirá a partir do sexto período (\autoref{sec:trajetoria_tecnologo}; \autoref{sec:trajetoria_engenharia}), mas não representa:

\begin{itemize}
    \item renúncia definitiva à outra graduação;
    \item transferência de curso;
    \item migração curricular (\autoref{sec:migracao_curricular});
    \item novo ingresso por processo seletivo;
    \item obtenção automática de dois diplomas;
    \item unificação dos dois cursos em uma única graduação; ou
    \item escolha irreversível que impeça a formação posterior na outra graduação.
\end{itemize}

O estudante concluirá a graduação escolhida. A possibilidade de prosseguir posteriormente na outra graduação somente existirá se for instituída por ato específico da UFU (\autoref{sec:permanencia}).

\textbf{Pendência institucional:} definir, antes da oferta, procedimento, prazo, órgão responsável, capacidade de cada curso e eventuais critérios acadêmicos para a escolha da graduação.

\section{Trajetória do curso superior de Tecnologia}
\label{sec:trajetoria_tecnologo}

Ao optar por concluir inicialmente o \ppccurso, o estudante cursa o sexto período — específico do \ppccurso, sem correspondência com a matriz da Engenharia de Automação, Robótica e Inteligência Artificial — e cumpre os demais requisitos próprios desta formação, previstos neste PPC: componentes curriculares obrigatórios e optativos, Atividades Curriculares de Extensão, Atividades Acadêmicas Complementares e Estágio Supervisionado Obrigatório (\autoref{sec:fluxo}; \autoref{sec:ace}; \autoref{sec:AAC}). A trajetória do \ppccurso é, portanto:

\begin{itemize}
    \item cinco períodos comuns da ABI (\autoref{sec:etapa_comum}) \textbf{mais} um período específico do \ppccurso, totalizando seis períodos — a duração mínima de \ppctempominimo já configurada para este curso (\autoref{chp:identificacao}), que esta organização em ABI não altera.
\end{itemize}

Concluídos os componentes e demais requisitos do sexto período, o estudante integraliza o \ppccurso e realiza a respectiva colação de grau (\autoref{sec:diplomas}), obtendo o diploma de \ppcgrau.

\section{Trajetória da Engenharia de Automação, Robótica e Inteligência Artificial}
\label{sec:trajetoria_engenharia}

Ao optar por concluir inicialmente a Engenharia de Automação, Robótica e Inteligência Artificial, o estudante passa a cursar os cinco períodos específicos dessa formação (sexto ao décimo período) e a cumprir os requisitos próprios previstos no respectivo PPC. A trajetória da Engenharia de Automação, Robótica e Inteligência Artificial é, assim:

\begin{itemize}
    \item cinco períodos comuns da ABI (\autoref{sec:etapa_comum}) \textbf{mais} cinco períodos específicos da Engenharia, totalizando dez períodos.
\end{itemize}

A duração de dez períodos pertence à Engenharia de Automação, Robótica e Inteligência Artificial e é apresentada aqui exclusivamente no contexto explicativo da organização da ABI — não é a duração do \ppccurso, nem substitui a duração de seis períodos registrada neste documento (\autoref{chp:identificacao}). O detalhamento da matriz curricular, dos componentes específicos, dos requisitos de integralização e do perfil de egresso da Engenharia de Automação, Robótica e Inteligência Artificial consta do respectivo PPC, elaborado separadamente.

\section{Conclusão da primeira graduação}
\label{sec:conclusao_primeira}

Qualquer que seja a graduação escolhida ao final do quinto período (\autoref{sec:escolha}), sua conclusão pressupõe o cumprimento integral dos requisitos previstos no respectivo PPC — não há dispensa ou abreviação de requisitos em razão da organização em ABI. Para o \ppccurso, esses requisitos são os descritos neste documento: componentes curriculares obrigatórios e optativos, carga horária mínima (\ppctempominimo), Atividades Curriculares de Extensão (\autoref{sec:ace}), Atividades Acadêmicas Complementares (\autoref{sec:AAC}) e Estágio Supervisionado Obrigatório (\autoref{sec:acc_conclusao}). Não há, neste curso, componente de Trabalho de Conclusão de Curso (\autoref{sec:acc_conclusao}).

\section{Permanência de vínculo}
\label{sec:permanencia}

As Normas Gerais da Graduação da UFU~\cite{ufu_congrad_46_2022} preveem permanência de vínculo para conclusão de outra habilitação, modalidade ou certificado de estudos oferecido pelo mesmo curso. Como esta proposta articula dois cursos com PPCs, códigos e diplomas distintos, essa norma não basta para autorizar a continuidade aqui descrita. Até a aprovação de ato institucional específico, não se assegura permanência de vínculo, ingresso na segunda graduação sem novo processo seletivo ou segundo diploma.

A eventual regulamentação deverá estabelecer que o mecanismo:

\begin{itemize}
    \item é uma possibilidade acadêmica, não uma etapa obrigatória da trajetória do estudante;
    \item constitui procedimento formal, mediante solicitação do próprio estudante — não é concedida de ofício;
    \item deve ser solicitada antes da colação de grau da primeira graduação concluída;
    \item está sujeita às normas acadêmicas institucionais vigentes, podendo depender da existência de vaga e de análise pelas instâncias competentes;
    \item permite que a continuidade na outra graduação ocorra sem caracterizar novo ingresso, migração curricular (\autoref{sec:migracao_curricular}) ou transferência de curso, observadas as normas institucionais aplicáveis a essa articulação entre dois cursos distintos; e
    \item não produz, por si só, o diploma da segunda graduação (\autoref{sec:continuidade_segunda}) — este depende do cumprimento integral dos requisitos curriculares ainda pendentes.
\end{itemize}

\textbf{Pendência institucional crítica:} definir a base normativa, a viabilidade administrativa, as vagas, o prazo, a forma do pedido, o órgão competente, os requisitos de deferimento, o registro acadêmico e os efeitos sobre colação de grau e diploma entre dois cursos distintos.

\section{Continuidade na segunda graduação}
\label{sec:continuidade_segunda}

Somente se houver regulamentação institucional e deferimento nos termos do futuro ato, o estudante poderá dar continuidade à outra graduação. Nessa hipótese:

\begin{itemize}
    \item os componentes curriculares comuns já integralizados são aproveitados, sem necessidade de nova integralização;
    \item aplicam-se as equivalências curriculares previamente definidas entre as duas matrizes (\autoref{sec:equivalencias_abi});
    \item o estudante deverá cursar os componentes específicos da segunda graduação ainda não integralizados e cumprir todos os demais requisitos curriculares dessa formação;
    \item o tempo necessário para a conclusão depende dos componentes e requisitos ainda pendentes, e não é automaticamente igual à duração da trajetória específica descrita nas Seções~\ref{sec:trajetoria_tecnologo} e~\ref{sec:trajetoria_engenharia} para quem a escolhe como primeira formação; e
    \item a conclusão da segunda graduação é seguida de nova colação de grau e da expedição do respectivo diploma (\autoref{sec:diplomas}).
\end{itemize}

A segunda graduação não é apresentada, em nenhuma hipótese, como diploma automático, habilitação adicional automática, apostilamento do primeiro diploma ou dupla diplomação automática — depende, em todos os casos, do cumprimento integral dos requisitos específicos ainda pendentes da respectiva formação.

\subsection{Do \ppccurso para a Engenharia de Automação, Robótica e Inteligência Artificial}

Na hipótese condicionada ao ato institucional, o estudante que concluir inicialmente o \ppccurso e obtiver autorização para a continuidade poderá prosseguir na Engenharia de Automação, Robótica e Inteligência Artificial, sujeito ao aproveitamento formal de estudos e aos requisitos da segunda formação. O fluxo pretendido é:

\begin{center}
Ingresso na ABI $\to$ cinco períodos comuns $\to$ escolha pelo \ppccurso $\to$ sexto período específico do \ppccurso $\to$ integralização do \ppccurso $\to$ permanência de vínculo (antes da colação de grau) $\to$ colação de grau no \ppccurso $\to$ continuidade na Engenharia $\to$ aproveitamento dos componentes comuns e das equivalências $\to$ componentes específicos restantes da Engenharia $\to$ integralização da Engenharia $\to$ segunda colação de grau $\to$ diploma da Engenharia de Automação, Robótica e Inteligência Artificial.
\end{center}

O conjunto de componentes específicos restantes corresponde, em princípio, aos cinco períodos específicos da Engenharia (sexto ao décimo período). O detalhamento exato desses componentes e de eventuais requisitos adicionais (estágio, trabalho de conclusão de curso e demais exigências próprias da Engenharia) consta da matriz curricular e da tabela de equivalências desse curso, elaboradas no âmbito do respectivo PPC.

\textcolor{red}{[CONFIRMAR A TABELA DE EQUIVALÊNCIAS ENTRE OS DOIS CURSOS]}

\subsection{Da Engenharia de Automação, Robótica e Inteligência Artificial para o \ppccurso}

Na hipótese condicionada ao ato institucional, o estudante que concluir inicialmente a Engenharia de Automação, Robótica e Inteligência Artificial e obtiver autorização para a continuidade poderá prosseguir no \ppccurso, sujeito ao aproveitamento formal de estudos. O fluxo pretendido é:

\begin{center}
Ingresso na ABI $\to$ cinco períodos comuns $\to$ escolha pela Engenharia $\to$ cinco períodos específicos da Engenharia $\to$ integralização da Engenharia $\to$ permanência de vínculo (antes da colação de grau) $\to$ colação de grau na Engenharia $\to$ continuidade no \ppccurso $\to$ aproveitamento dos componentes comuns e das equivalências $\to$ requisitos específicos restantes do \ppccurso $\to$ integralização do \ppccurso $\to$ segunda colação de grau $\to$ diploma do \ppccurso.
\end{center}

Como os cinco primeiros períodos são compartilhados entre as duas matrizes (\autoref{sec:etapa_comum}), e os componentes do sexto período específico do \ppccurso (\autoref{sec:trajetoria_tecnologo}) têm como pré-requisitos apenas carga horária acumulada e correquisitos entre si — sem dependência de nenhum componente posterior ao quinto período —, é razoável esperar que os requisitos residuais, neste sentido de continuidade, correspondam ao sexto período específico do \ppccurso (Empreendedorismo e Inovação, Atividades Curriculares de Extensão III, Proteção e Acionamentos Industriais, Projeto Interdisciplinar, Instalações Elétricas e sua atividade experimental, disciplina optativa e Estágio Supervisionado Obrigatório — \autoref{sec:fluxo}). Essa correspondência, no entanto, depende da confirmação de que os componentes específicos da Engenharia (sexto ao décimo período) não introduzem exigências próprias que também precisem ser satisfeitas antes do reconhecimento como equivalentes ao quinto período do \ppccurso.

\textcolor{red}{[CONFIRMAR OS COMPONENTES E REQUISITOS RESIDUAIS PARA O ESTUDANTE QUE CONCLUIR PRIMEIRO A ENGENHARIA E DEPOIS SOLICITAR A CONTINUIDADE NO TECNÓLOGO]}

\section{Equivalências e aproveitamento de estudos}
\label{sec:equivalencias_abi}

O reconhecimento de componentes curriculares entre o \ppccurso e a Engenharia de Automação, Robótica e Inteligência Artificial, para os fins descritos neste capítulo, apoia-se em dois conceitos distintos:

\begin{description}
    \item[Equivalência curricular] Reconhecimento formal de compatibilidade entre componentes curriculares das duas matrizes, considerando conteúdo, carga horária, competências e resultados de aprendizagem, entre outros critérios institucionais.
    \item[Aproveitamento de estudos] Procedimento acadêmico pelo qual componentes curriculares já integralizados são reconhecidos com base nas equivalências previamente aprovadas e nas normas institucionais vigentes, dispensando o estudante de cursá-los novamente.
\end{description}

Como os cinco primeiros períodos são comuns às duas matrizes (\autoref{sec:etapa_comum}), o aproveitamento dos componentes dessa etapa, ao se solicitar a permanência de vínculo, decorre diretamente dessa identidade curricular. O aproveitamento não constitui dispensa genérica ou automática de todos os requisitos da segunda graduação: os componentes específicos ainda não integralizados, e os demais requisitos próprios da formação (estágio, atividades complementares, atividades de extensão e demais exigências aplicáveis), permanecem exigíveis (\autoref{sec:continuidade_segunda}).

\textcolor{red}{[CONFIRMAR A TABELA DE EQUIVALÊNCIAS ENTRE OS DOIS CURSOS]}

\section{Colações de grau e expedição dos diplomas}
\label{sec:diplomas}

A conclusão de cada graduação integrante da ABI é seguida de colação de grau e de expedição de diploma próprios e distintos:

\begin{itemize}
    \item a conclusão do \ppccurso confere o diploma de \ppcgrau em Automação Industrial, mediante colação de grau específica deste curso;
    \item a conclusão da Engenharia de Automação, Robótica e Inteligência Artificial confere o diploma de Bacharel nessa formação, mediante colação de grau específica desse curso, nos termos do respectivo PPC.
\end{itemize}

Na hipótese de a continuidade ser regulamentada e de o estudante concluir as duas graduações, haverá colação de grau e diploma próprios para cada curso, conforme o futuro ato institucional e os procedimentos de registro acadêmico. Não há emissão automática ou simultânea dos dois diplomas, nem diploma único que os represente conjuntamente.

\textcolor{red}{[CONFIRMAR PROCEDIMENTO PARA A SEGUNDA COLAÇÃO DE GRAU]}

\textcolor{red}{[CONFIRMAR O REGISTRO ACADÊMICO DA SEGUNDA GRADUAÇÃO]}

\section{Implantação e inexistência de migração curricular}
\label{sec:migracao_curricular}

O \ppccurso é curso novo — encontra-se, no momento da elaboração deste PPC, em fase de proposta (\autoref{chp:justificativa}) —, e a Engenharia de Automação, Robótica e Inteligência Artificial, igualmente, é proposta de curso novo. A organização em ABI descrita neste capítulo é, portanto, a forma de ingresso e de articulação entre os dois cursos desde sua implantação, e não uma migração de estudantes de uma matriz curricular anterior para a estrutura aqui descrita — distinção relevante entre dois conceitos que não devem ser confundidos:

\begin{description}
    \item[Migração curricular] Mudança de estudantes vinculados a um currículo anterior de um mesmo curso para um currículo novo desse curso — situação que não se aplica a nenhum dos dois cursos integrantes desta ABI, por serem ambos cursos novos, sem currículo anterior ao aqui descrito.
    \item[Continuidade proposta entre cursos] Mecanismo ainda dependente de ato institucional específico, distinto da migração curricular e da permanência de vínculo já disciplinada para habilitações, modalidades ou certificados de um mesmo curso (\autoref{sec:permanencia}).
\end{description}

Não há migração ou transferência automática entre os cursos. A etapa curricular comum não cria, por si só, direito de ingresso na outra graduação; qualquer continuidade dependerá da regulamentação indicada na \autoref{sec:permanencia}.
